\section{Early Reasoning 文献：RLP 与 Quiet-STaR、RPT、RLPT 的差别}

让模型“在说话前思考”并非 RLP 首创。Quiet-STaR、Reinforcement Pre-Training（RPT）与 Reinforcement Learning on Pre-Training Data（RLPT）都把 reasoning 提前。真正需要比较的是 reward 从哪里来、在什么粒度计算、是否需要外部 verifier、以及 reasoning trace 是否直接获得 credit。

\subsection{从 Quiet-STaR 到 reinforcement pretraining}

文献页把 Quiet-STaR、RPT、RLPT 与 RLP 放在同一时间线上。Quiet-STaR 在普通文本中生成并评估 thought；RPT 把下一 token 正确性转成可验证的 RL reward；RLPT 使用 pretraining data 构造更长片段的强化学习任务；RLP 则用 thought 对真实下一 token 的 information gain 作为 dense intrinsic signal。

\lecturefigure{slide-038.jpg}{Early-reasoning 路线：Quiet-STaR、RPT、RLPT 与 RLP}{Stanford 官方录像，00:42:20--00:43:13}

\begin{knowledgebox}{术语表：相邻方法解决什么}
\begin{tabularx}{\textwidth}{L{0.18\textwidth}X X}
\toprule
方法 & 核心机制 & 与 RLP 的关键差别 \\
\midrule
NTP & 直接最大化真实下一 token likelihood & 无显式 thought action / reward \\
Quiet-STaR & 在文本中采样 thought 并衡量其效用 & 早期 self-thought 路线，训练细节与 RLP 不同 \\
RPT & 把 next-token reasoning 变成可验证 RL 任务 & 使用筛选与 sparse binary correctness reward \\
RLPT & 在 pretraining data 上做 segment-level RL & reward 粒度和 verifier 设计不同 \\
RLP & 用 thought 带来的 log-likelihood gain 奖励 rollout & dense、token-level、no-think counterfactual \\
\bottomrule
\end{tabularx}
\end{knowledgebox}

\subsection{三条比较轴：reward source、granularity、emergence}

NTP 没有单独 reward；RPT/RLPT 使用外部或可验证信号，通常较 sparse；RLP 的 reward 来自内部 reasoned predictor 与 EMA baseline 的差，并在 token 位置上连续取值。课程把 RLP 标为 explicit/strong reasoning emergence，是作者对训练信号的概括，不是对内部认知过程的直接测量。

\lecturefigure{slide-039.jpg}{NTP、RPT/RLPT 与 RLP 的 reward 与 reasoning 对照}{Stanford 官方录像，00:43:13--00:43:56}

\begin{warningbox}{Verifier-free 的精确定义}
RLP 不需要另一个答案判定模型或人工标注 verifier，但它仍以语料中的真实下一 token 为监督，并使用 EMA baseline。因而更准确的说法是“external-verifier-free”，而不是“没有监督或没有参照”。
\end{warningbox}

\subsection{Matched RPT comparison：机制与数字一起读}

定性表只能说明设计差异，不能说明哪个方法在相近预算下更有效；因此课程接着用 matched setting 把 reward mechanism 与量化结果放在同一页。课程在 Qwen3-1.7B-Base、OmniMath、170M token 等相近设置下比较 RPT 与 RLP，报告 RLP 平均约高 4\%。机制解释是：RPT 先用外部 filter 找可用位置，再给 sparse binary reward；RLP 可在任意选定 token 上计算 dense information gain，并把 credit 连到 reasoning trace。

\lecturefigure{slide-040.jpg}{RLP 与 RPT 的 matched-setting 比较}{Stanford 官方录像，00:43:56--00:45:56}

\begin{knowledgebox}{读图：哪些变量仍可能影响比较}
即使模型、数据与 token 数接近，prompt、rollout 数、thought length、optimizer、KL/clipping、sample filtering 和 compute 仍可能不同。图支持“在作者复现实验里 RLP 更好”，不等于对所有 RPT 实现的理论支配。
\end{knowledgebox}

\subsection{本章小结}

RLP 属于 early-reasoning / reinforcement-pretraining 家族，其辨识度来自 no-think counterfactual 与 dense information-gain reward。比较相邻方法时必须同时看 reward source、granularity、筛选成本与 compute，而不是只比较方法名称或单个平均分。

